\paragraph*{04.05.07.05 Finanzen überwachen, um Abweichungen vom Projektplan zu identifizieren und zu korrigieren}
\kcilabel{para:04.05.07.05_FinanzenÜberwachen}{04.05.07.05}
\label{para:04.05.07.05_FinanzenÜberwachen}

% Task
\par  Die in \ref{para:04.05.07.01_Projektkosten} beschriebenen Kostenschätzungen und in einen Kostenplan überführten Werte bildeten die Grundlage für die finanzielle Steuerung des Projekts. Meine Aufgabe bestand darin, die tatsächlichen Kosten der Entwicklung und Implementierung des \gls{wsr} und \gls{pb} kontinuierlich zu überwachen, um Abweichungen vom genehmigten Kostenplan frühzeitig zu erkennen und geeignete Korrekturmaßnahmen einzuleiten.

% Action
\par Hierzu entwickelte ich gemeinsam mit der \gls{pl} ein monatliches Kostenreporting mit wellenspezifischem Soll-Ist-Vergleich und Forecast. Dafür strukturierte ich die Entwicklungsaufwände der aktuellen Umsetzungswelle, stimmte die Kostenzuordnung mit den verantwortlichen Liefer- und Fachbereichen ab und verdichtete die Ergebnisse zu einer belastbaren Entscheidungsgrundlage für die Steuerung. Erkannte Abweichungen wurden von mir hinsichtlich Ursache, Auswirkung und Handlungsbedarf bewertet. Ein konkreter Fall war eine Forecast-Abweichung in Welle 1.1 aufgrund zusätzlicher Integrationsaufwände im Umfang von 5 PT; ich bereitete die Entscheidungsoptionen transparent auf (Leistungsanpassung, Terminverschiebung, zusätzliche Mittel) und führte die Entscheidung in das zuständige Gremium über. Bei drohenden Überschreitungen leitete ich gemeinsam mit den Beteiligten geeignete Maßnahmen ab, etwa die Priorisierung von Anforderungen, die zeitliche Verschiebung nachgelagerter Funktionalitäten oder die Nachschärfung von Leistungszuschnitten. Die Berichterstattung erfolgte in den etablierten Entscheidungs- und Berichtsstrukturen gemäß \ref{para:04.03.02.05_EntscheidungsUndBerichtsstrukturenUndQualitätsanforderungen} und \ref{para:04.04.05.05_EntscheidungenTreffenDurchsetzenUndÜberprüfen}; bei Abweichungen oberhalb des definierten Schwellenwerts von 10\% eskalierte ich die Entscheidungsvorlage fristgerecht.

% Result
\par Durch das etablierte Kostenreporting konnten finanzielle Abweichungen frühzeitig erkannt und adressiert werden. Die Prognosegüte verbesserte sich von 25\% auf 10\%, und die mittlere Reaktionszeit vom Erkennen einer Abweichung bis zur Beschlussfassung sank auf 3 Arbeitstage. Dies ermöglichte eine belastbare Steuerung der Entwicklungsbudgets je Welle und trug dazu bei, den genehmigten Kostenrahmen trotz der hohen fachlichen und technischen Dynamik im Programm einzuhalten. Gleichzeitig erhöhte die Transparenz über Kostentreiber, Forecast-Entwicklungen und erforderliche Korrekturmaßnahmen die Entscheidungssicherheit für die weitere Umsetzungsplanung.

% Reflect
\par Für komplexe Transformationsvorhaben reicht eine einmalige Kostenschätzung nicht aus; entscheidend ist ein kontinuierliches Finanzmonitoring mit klarer Zuordnung von Kostentreibern, belastbaren Forecasts und definierten Reaktionsmechanismen. Ich habe daraus mitgenommen, dass insbesondere bei hoher Scope-Dynamik und wellenweiser Umsetzung finanzielle Abweichungen nur dann wirksam beherrscht werden können, wenn Kostensteuerung eng mit Leistungsumfang, Priorisierung und Entscheidungswegen verzahnt ist. Für künftige Vorhaben würde ich die Schwellenwerte und Eskalationspfade bereits zu Projektbeginn verbindlich festlegen und die Wirksamkeit der Gegenmaßnahmen anhand weniger, klar definierter Steuerungskennzahlen regelmäßig nachhalten.